% Chapter 1 — Visión general del proyecto
\chapter{Visión General del Proyecto}

\section{Descripción}

\textbf{cv-lit} es un sistema de detección automática de línea de costa a partir de imágenes
oblicuas capturadas por \textbf{6 cámaras fijas} instaladas en el frente litoral de
\textbf{Guardamar del Segura} (Alicante, España).

\section{Cliente y colaboradores}

\begin{description}
  \item[Cliente] Ayuntamiento de Guardamar del Segura.
  \item[Colaboración técnica] Instituto de Ecología Litoral (IEL) — suministra
        los GCPs GNSS para la calibración de homografía.
  \item[Equipo de desarrollo] 1 ingeniero a tiempo parcial, supervisado por la
        Universidad de Alicante.
  \item[Duración] 6 meses.
\end{description}

\section{Objetivo}

Monitorización automática y continua de la playa seca produciendo como
producto final un fichero \textbf{GeoJSON proyectado en EPSG:25830}
(ETRS89 / UTM zona 30N).

\begin{table}[h!]
\centering
\begin{tabular}{lp{9cm}}
\toprule
\textbf{Campo GeoJSON} & \textbf{Descripción} \\
\midrule
\texttt{ID\_Camara}     & Identificador de la cámara (1--6) \\
\texttt{Timestamp}      & Fecha y hora UTC de la imagen \\
\texttt{Confianza\_IA}  & Score de confianza del modelo SAM (0--1) \\
\texttt{Area\_Seca\_m2} & Superficie de playa seca medida en m² \\
\bottomrule
\end{tabular}
\caption{Atributos obligatorios del GeoJSON de salida}
\end{table}

\section{Restricciones no negociables}

\begin{itemize}
  \item Toda implementación \textbf{debe} producir GeoJSON proyectado en \textbf{EPSG:25830}.
  \item La fase horaria prioritaria de adquisición es \textbf{12:00h solar local}.
  \item Modelo de segmentación: \textbf{SAM} (Segment Anything Model, Meta AI).
  \item Cada cámara tiene su propio perfil de calibración independiente.
\end{itemize}

\section{Entregables finales}

\begin{enumerate}
  \item Aplicación ejecutable configurada y documentada.
  \item Perfiles de calibración por cámara (6 ficheros \texttt{.npy}/\texttt{.json}).
  \item Manual breve de uso.
  \item Informe inicial de validación con métricas RMSE y MAE.
\end{enumerate}
